\documentclass{ximera}
\title{Which Tanks Empty Faster?}

\newcommand{\pskip}{\vskip 0.1 in}

\begin{document}
\begin{abstract}
Draining tanks.
\end{abstract}
\maketitle



\section{Preliminaries}

\section{Thoughts on Differential Equation in Math 152}

It is possible to integrate (ha!) differential equations in a meaningful way rather than saving it for the end.

For me, this approach starts with changing the focus from using definite integrals and Part 2 (?) of the Fundamental Theorem to compute \emph{numbers} to using definite integrals and Part 1 of the FT to find expressions for functions. It makes generalizing the point-slope equation of a line a fundamental part of the course.

Here are two examples.

\begin{example} \label{Exkr340}

\begin{enumerate}

\item \emph{Math 99} During a $50$-mile trip a car gets a constant gase mileage of $40$ miles/gal. When the trip odometer reads $32$ miles, the car has $9$ gallons of gas in its tank.

Find a function
\[
       G = h(s) \, , \, 0\leq s \leq 50 ,       
\]
that expresses the number of gallons of gas in the tank in terms of the trip odometer reading (in miles).


\item \emph{Math 152} The function 
\[
    m = f(s) = ??? \, , \, 0\leq s \leq 50 ,
\]
expresses the gas mileage of a car (measured in miles/gal) in terms of the trip odometer reading (measured in miles) during a $50$-mile trip. When the trip odometer reads $32$ miles, the car has $9$ gallons of gas in its tank.

Find a function
\[
     G = h(s) \, , \, 0\leq s \leq 50 ,
\]
that expresses the number of gallons of gas in the tank in terms of the trip odometer reading (in miles).

\end{enumerate}

\begin{explanation} 
\begin{enumerate}
\item 
\[
    G = h(s) = 9 - \frac{1}{40} \left( s-32  \right) \, , \, 0\leq s \leq 50. 
\]

\item Two approaches. One starts by adding the differential number of gallons of gas
\[
    dG = - \frac{ds}{f(s)}
\]
the car burns over the intervals $s \in [u, u+du]$ miles as the trip odometer ranges from $s=32$ miles to $s=u$ miles. 

The other starts with an initial value problem
\[
   \frac{dG}{ds} = - \frac{1}{f(s)} \, , \, G=9 \text{ when } s=32.
\]

Either way we get
\[
  G = h(s) = 9 - \int_{32}^s \frac{du}{f(u)} \, , \, 0\leq s \leq 50. 
\]

The two approaches are logically equivalent if we use definite integrals to solve the differential equation. Using indefinite integrals misses the point. The integral stops being a sum and the computation loses meaning.

\end{enumerate} 
\end{explanation}
\end{example}


Here's a more interesting example.

\section{Which Tanks Drain Faster?}
Fill a tank with water, punch a small hole in the bottom of the tank, and let the water drain out. The problem is to find a function that expresses the depth of water in terms of time.

\begin{onlineOnly}
    \begin{center}
\desmos{zltxdekv7z}{900}{600}
\end{center}
\end{onlineOnly}

\href{https://www.desmos.com/calculator/zltxdekv7z}{152: Draining Cylinders}

The given information:

\begin{enumerate}
\item The cross-sectional area $A=f(h)$ of the tank as a function of depth.

\item The area $a$ of the small hole.

\item The assumption (or theoretical fact) that the water exits the hole with speed
\[
  v = g(h) =\sqrt{2gh},
\] 
where $h$ is the depth of the water and $g$ the magnitude of the graviational acceleration near the earth's surface.
\end{enumerate}



\section{A Hemispherical Tank}

\begin{onlineOnly}
    \begin{center}
\desmos{ggqodhz3jv}{900}{600}
\end{center}
\end{onlineOnly}

\href{https://www.desmos.com/calculator/ggqodhz3jv}{152: Draining Cylinders 2}



\end{document}
